% chapters/07/section-06.tex
\section{Menganalisis Tabrakan yang Salah}
\label{sec:analisis-tabrakan}

Penanganan tabrakan yang salah hampir selalu berjalan tanpa pesan kesalahan.
Program tetap berjalan, gambar tetap muncul, dan hanya perilakunya yang janggal.
Bagian ini melatih cara menemukan dan memperbaikinya, sesuai indikator \ind{c}
pada bab ini.

\subsection{Cara Menelusuri}

Empat langkah berikut berlaku umum untuk seluruh kesalahan yang tidak
menimbulkan pesan kesalahan, bukan hanya untuk tabrakan.

\begin{langkah}
  \item \textbf{Dapatkan gejala yang dapat diulang.} Tentukan langkah pasti yang
        memunculkan gejala itu: dari posisi mana, bergerak ke arah mana, dengan
        kecepatan berapa. Gejala yang tidak dapat diulang tidak dapat
        diperbaiki.
  \item \textbf{Persempit wilayahnya.} Tentukan bagian program yang terlibat.
        Apabila gejala hanya muncul ketika menyentuh bongkahan, maka
        pemeriksaan bongkahanlah yang perlu diperiksa lebih dahulu.
  \item \textbf{Pasang alat ukur.} Tampilkan nilai yang Anda duga menjadi
        penyebabnya: posisi sebelum dan sesudah pembaruan, nilai \code{rect},
        dan kecepatan pada kedua sumbu. Pada program bab ini, penghitung
        tumpang tindih memainkan peran itu.
  \item \textbf{Ubah satu hal pada satu waktu.} Ubah dugaan penyebabnya, jalankan
        kembali, dan bandingkan. Mengubah dua hal sekaligus membuat Anda tidak
        tahu mana yang berpengaruh.
\end{langkah}

\subsection{Tiga Gejala yang Paling Sering}

Tabel~\ref{tab:gejala-tabrakan} memuat tiga gejala yang paling sering muncul
beserta penyebab dan perbaikannya.

\begin{table}[htbp]
\centering
\caption{Gejala tabrakan yang salah, penyebab, dan perbaikannya}
\label{tab:gejala-tabrakan}
{\setlength{\tabcolsep}{4pt}\footnotesize%
\begin{tabularx}{\textwidth}{@{}L{3.4cm} L{4.2cm} Y@{}}
\toprule
\headrow \textbf{Gejala} & \textbf{Penyebab} & \textbf{Perbaikan} \\
\midrule
Benda menembus bongkahan yang tipis, terutama ketika bergerak cepat &
Pada satu bingkai, perpindahan benda melampaui tebal bongkahan, sehingga
keduanya tidak pernah beririsan pada saat pemeriksaan. Gejala ini disebut
\istilah{tunneling}. &
Batasi panjang langkah, bagi perpindahan menjadi beberapa langkah kecil, atau
tebalkan bongkahan. Perhitungannya ada pada
Bagian~\ref{sec:contoh-bab-7}. \\
\addlinespace
Benda yang menabrak sisi bongkahan tiba-tiba terangkat ke atasnya &
Kedua sumbu diperiksa sekaligus, sehingga program tidak mengetahui sisi mana
yang ditabrak dan selalu memilih penyelesaian pada sumbu tegak. &
Pisahkan pemeriksaan menjadi dua tahap seperti pada
Bagian~\ref{sec:per-sumbu}. \\
\addlinespace
Benda bergetar halus di permukaan bongkahan &
Kecepatan tegak tidak dinolkan ketika tabrakan diselesaikan, sehingga benda
terus tertekan ke dalam bongkahan dan posisinya diperbaiki setiap bingkai. &
Nolkan kecepatan pada sumbu yang diselesaikan. \\
\addlinespace
Benda dinilai bersentuhan padahal gambarnya masih berjarak &
Benda berbentuk miring atau berlubang, sehingga persegi pembatasnya jauh lebih
besar daripada gambarnya. &
Pakai \istilah{mask} sesuai pertimbangan pada
Tabel~\ref{tab:pilihan-metode}. \\
\bottomrule
\end{tabularx}}
\end{table}

\subsection{Kebiasaan yang Mencegahnya}

Empat kebiasaan berikut mencegah sebagian besar kesalahan di atas muncul sejak
awal.

\begin{itemize}
  \item \textbf{Selesaikan satu sumbu sampai tuntas sebelum menyentuh sumbu
        berikutnya.} Ini kebiasaan yang paling banyak mencegah kesalahan.
  \item \textbf{Tampilkan keadaan yang Anda curigai pada layar atau terminal.}
        Angka yang terlihat akan menjawab pertanyaan yang tidak dapat dijawab
        oleh gambar.
  \item \textbf{Sediakan pemeriksa mandiri di dalam program}, sebagaimana
        penghitung tumpang tindih pada bab ini. Pemeriksa itu menjawab
        pertanyaan ``apakah keadaan yang terlarang sedang terjadi?'' pada
        setiap bingkai, bukan hanya ketika Anda memperhatikannya.
  \item \textbf{Jalankan program pada beberapa kecepatan.} Kesalahan
        \istilah{tunneling} dan kesalahan pembulatan hanya muncul pada
        kecepatan tertentu, sehingga pengujian pada satu kecepatan saja tidak
        mencukupi.
\end{itemize}

\begin{catatabox}
Ketiga gejala pada Tabel~\ref{tab:gejala-tabrakan} memiliki
satu sifat yang sama: semuanya merupakan akibat dari pilihan cara, bukan dari
kesalahan penulisan. Tidak ada satu pun yang dapat ditemukan oleh pemeriksa
sintaksis. Karena itu, kemampuan menganalisis seperti pada bagian ini merupakan
kemampuan yang membedakan program yang ``berjalan'' dari program yang
``benar''.
\end{catatabox}

\subsection{Mengakhiri Blok I}

Bab ini menutup \CPMK{1}, yaitu kemampuan merancang \istilah{game loop} dan
logika fisika sederhana. Dengan demikian, seluruh perkakas Blok~I telah
lengkap.

\begin{table}[htbp]
\centering
\caption{Perkakas Blok I yang telah lengkap beserta bab yang membahasnya}
\label{tab:perkakas-blok-1}
{\setlength{\tabcolsep}{4pt}\footnotesize%
\begin{tabularx}{\textwidth}{@{}L{3.6cm} L{3.0cm} Y@{}}
\toprule
\headrow \textbf{Perkakas} & \textbf{Dibahas pada} & \textbf{Kegunaan} \\
\midrule
Lingkungan kerja dan repositori & Bab~III &
Menjalankan program dan menyimpan riwayat pekerjaan. \\
\addlinespace
Siklus permainan dan pengendalian waktu & Bab~IV &
Menjalankan permainan terus-menerus dengan laju yang terkendali. \\
\addlinespace
Vektor, masukan, dan gerak berbasis waktu & Bab~V &
Menggerakkan benda dengan kecepatan yang sama pada setiap komputer. \\
\addlinespace
Gravitasi, lompatan, dan pantulan & Bab~VI &
Memberi benda sifat jatuh, melompat, dan memantul. \\
\addlinespace
Kelas \code{Sprite}, wadah \code{Group}, dan tabrakan & Bab~VII &
Mengelola banyak benda sekaligus dan menangani persentuhannya. \\
\bottomrule
\end{tabularx}}
\end{table}

Bab~VIII sampai Bab~XI akan memakai seluruh perkakas itu untuk keperluan baru:
membuat benda yang mengambil keputusan sendiri.
